\documentclass[11pt]{article}

\usepackage[margin=1in]{geometry}
\usepackage[T1]{fontenc}
\usepackage[utf8]{inputenc}
\usepackage{lmodern}
\usepackage{microtype}
\usepackage{amsmath,amssymb,amsthm,mathtools}
\usepackage{booktabs}
\usepackage{array}
\usepackage{enumitem}
\usepackage{listings}
\usepackage{xurl}
\usepackage[hidelinks]{hyperref}

\newtheorem{definition}{Definition}
\newtheorem{lemma}{Lemma}
\newtheorem{theorem}{Theorem}
\newtheorem{remark}{Remark}

\newcommand{\TV}{\mathrm{TV}}
\newcommand{\KL}{\mathrm{KL}}
\newcommand{\eps}{\varepsilon}
\newcommand{\logtwo}{\log_2}

\lstset{basicstyle=\footnotesize\ttfamily,breaklines=true,frame=single,columns=fullflexible}

\title{MC-EQ and CoverSketch:\\
Destination Privacy via Markov-Chain Equalization in Anonymous DHT Lookups}
\author{}
\date{Draft v0.46 (2026-02-28; archive v125)}

\begin{document}
\maketitle

\begin{abstract}
Destination privacy is the promise that an observer cannot reliably infer \emph{which key} a user is querying.
In anonymous overlays, this typically reduces to ensuring that routing transcripts are indistinguishable across destinations.
This paper introduces a compact equalization interface for destination privacy: \emph{MC-EQ} (Markov-Chain Equalization) and its deployable ``cover set'' implementation pattern, \emph{CoverSketch}.

MC-EQ treats routing as a (possibly history-dependent) Markov process over public states (committees, buckets, or routing regions).
A mechanism is destination-private if, at each step, its conditional next-state distribution stays within a declared total-variation budget of a public \emph{reference kernel}.
CoverSketch realizes this by restricting each state's outgoing support to a public cover set and equalizing probabilities over that set.

We keep the paper brief by importing the full AnonDHT conditioning model and MUCC plumbing from the state series \cite{series:state2}.
Our contribution is the equalization interface, simple composable leakage bounds (TV, ratio/$D_\infty$, and session scaling $T\log_2(1+2k\delta_{eq})$ under uniform covers), and a receipt-friendly parameterization suitable for the release workflow \cite{series:release}.
\end{abstract}

\paragraph{Units.}
Endpoint budgets in this archive are published in \emph{bits} (see the notation/model sheet, Synthesis~8).\cite{series:notation}
Accordingly, any conversion from a per-step $D_\infty$ cap to an odds-inflation or MaxL budget is rendered with $\log_2$.



\paragraph{Scope.}
We define a compact destination-privacy interface for Anonymous-DHT routing transcripts: Markov-Chain Equalization (MC-EQ) and its cover-set implementation pattern (CoverSketch).

\paragraph{Imports.}
We import the AnonDHT transcript/state conditioning model (including MUCC) from the state series.\cite{series:state2}
We treat endpoint conversions (TV $\leftrightarrow$ odds/max-leakage sizing) and release receipt semantics as owned by the companion papers.\cite{series:endpointbridge,series:release}

\paragraph{Exports.}
The MC-EQ contract (reference kernel + per-step TV budgets), a deployable CoverSketch construction, and simple composable leakage bounds exportable as receipt parameters.

\paragraph{Artifact policy.}
Source-only; supporting simulations and kernel-estimation scripts are available on request.\cite{series:artifactpolicy}



\section{MC-EQ: the equalization contract}
Let $S$ be a finite set of public routing states (e.g., committees).
A lookup for destination $d$ produces a (random) state sequence $X_0,\dots,X_T$.
Write the (possibly history-dependent) transition rule as
\[
P_d(x_{t+1}\mid h_t) \quad\text{where } h_t=(x_0,\dots,x_t,\text{public side-info}).
\]

\begin{definition}[Reference kernel]
A \emph{reference kernel} $K(\cdot\mid x)$ is a public Markov kernel on $S$.
It represents the ``idealized'' destination-independent next-state distribution.
\end{definition}

\begin{definition}[MC-EQ($\eps_t$)]
A mechanism satisfies \textbf{MC-EQ} with parameters $(\eps_t)_{t=0}^{T-1}$ and reference kernel $K$ if for every destination $d$, every time $t$, and every history $h_t$ with current state $x_t$,
\[
\TV\big(P_d(\cdot\mid h_t),\,K(\cdot\mid x_t)\big)\;\le\; \eps_t.
\]
\end{definition}

The contract is local-in-time and history-robust.
It is designed for certification: $K$ and $(\eps_t)$ are simple objects that can be published and verified.

\subsection{Why TV?}
TV is operational: it upper-bounds an observer's single-shot advantage in distinguishing two distributions.
It composes cleanly across steps via standard couplings.
For other leakage measures, we can convert TV bounds under mild ``floor'' assumptions (\S\ref{sec:linfty}).

\section{Composed transcript bounds}
The main question: what does local MC-EQ buy you for the whole transcript?

\subsection{TV composition (coupling view)}
\begin{theorem}[TV bound for path distributions]
Suppose a mechanism satisfies MC-EQ($\eps_t$) with reference kernel $K$.
Let $\mathcal{L}_d$ be the law of the length-$T$ transcript $(X_0,\dots,X_T)$ for destination $d$.
Then for any two destinations $d,d'$,
\[
\TV(\mathcal{L}_d,\mathcal{L}_{d'})\;\le\; \sum_{t=0}^{T-1} \eps_t.
\]
\end{theorem}
\begin{proof}[Proof sketch]
Couple the two processes stepwise.
At each step, the conditional next-state laws for $d$ and $d'$ are each within $\eps_t$ of the same reference kernel $K(\cdot\mid x_t)$.
A maximal coupling shows that the probability the coupled chains diverge at step $t$ is at most $\eps_t$.
Union bound over steps yields the claim.
\end{proof}

This TV bound is intentionally simple: it is conservative but receipt-friendly.
Sharper bounds (e.g., using spectral contraction when $K$ mixes) can be imported from \texttt{[[2026.01.23 - Mathematics: Spectral Anonymity and Optimal Distinguishing Bounds for Random-Walk Delegation in (Possibly Directed) Overlays]]}.

\subsection{$D_\infty$ / ratio conversion and session scaling}
\label{sec:linfty}
TV bounds alone do not control likelihood ratios if the reference kernel assigns tiny mass.
We therefore make a standard ``mass floor'' assumption.

\begin{definition}[Floor / mass minimum]
For a kernel $K$, define
\[
q_{\min}\;=\;\min\{K(y\mid x):\; x\in S,\; y\in \mathrm{supp}(K(\cdot\mid x))\}.
\]
\end{definition}
If $K(\cdot\mid x)$ is uniform over a cover set of size $k(x)$, then $q_{\min}=\min_x 1/k(x)$.

\begin{lemma}[TV to pointwise ratio under a floor]
\label{lem:tv-to-ratio}
Let $Q$ be a distribution with support $S$ and minimum mass $q_{\min}:=\min_{y\in S}Q(y)>0$.
If $\TV(P,Q)\le \delta$, then for all $y\in S$,
\[
Q(y)-2\delta \le P(y) \le Q(y)+2\delta.
\]
Consequently, for all $y\in S$,
\[
\frac{P(y)}{Q(y)}\le 1+\frac{2\delta}{q_{\min}},
\]
and if $\delta<q_{\min}/2$ then also
\[
\frac{P(y)}{Q(y)}\ge 1-\frac{2\delta}{q_{\min}}.
\]
\end{lemma}

\begin{proof}
Since $\sum_{z}|P(z)-Q(z)|=2\TV(P,Q)\le 2\delta$, we have $|P(y)-Q(y)|\le 2\delta$ for every $y$, yielding the additive bounds.
Dividing by $Q(y)\ge q_{\min}$ gives the ratio bounds.
\end{proof}

\paragraph{Per-step ratio and $D_\infty$.}
Applying Lemma~\ref{lem:tv-to-ratio} with $Q(\cdot)=K(\cdot\mid x)$ yields a per-step likelihood-ratio cap:
\[
D_\infty\!\big(P_d(\cdot\mid h_t)\,\|\,K(\cdot\mid x_t)\big)\;\le\;\logtwo\!\Big(1+\frac{2\eps_t}{q_{\min}}\Big).
\]
In discrete settings, maximal leakage admits simple upper bounds in terms of worst-case likelihood ratios and related max-divergences (see, e.g., \cite{issa_maxleak}).
In this series, endpoint conversions and odds-factor interpretations are centralized in the endpoint-bridge note.\cite{series:endpointbridge}

\paragraph{Session scaling (uniform support case).}
A reviewer- and operator-friendly knob is how leakage scales with the number of observable steps $T$.
If the reference kernel is uniform over a cover set of size at most $k$ for every state, then $q_{\min}=1/k$ and a single-step MC-EQ slack $\delta_{eq}$ yields the familiar factor $1+2k\delta_{eq}$.
Instantiating the series' generic pointwise-to-session export for per-history likelihood-ratio caps (owned by the scheduling/compiler backbone) gives the additive session bound:
\[
\mathcal{L}_{\max}(C\to O_{1:T}) \;\le\; T\,\log_2\!\Big(1+2k\,\delta_{eq}\Big).
\]
(For the generic ``per-step ratio $\Rightarrow$ session odds'' lemma used throughout the archive, see the scheduling/compiler note.\cite{series:schedcompiler})

\begin{table}[t]
\centering
\small
\begin{tabular}{@{}r r r@{}}
\toprule
$k$ & $\delta_{eq}$ & bound at $T=50$ (bits) \\
\midrule
8 & 0.005 & 5.55 \\
8 & 0.010 & 10.71 \\
8 & 0.020 & 20.03 \\
16 & 0.005 & 10.71 \\
16 & 0.010 & 20.03 \\
16 & 0.020 & 35.68 \\
32 & 0.005 & 20.03 \\
32 & 0.010 & 35.68 \\
32 & 0.020 & 59.45 \\
\bottomrule
\end{tabular}
\caption{Analytic maximal-leakage upper bound $T\log_2(1+2k\delta_{eq})$ for representative parameters (uniform $k$-support reference).}
\label{tab:session-bound}
\end{table}

\begin{remark}[Sharper bounds via mixing]
The bound above is deliberately conservative and receipt-friendly.
If the reference chain mixes rapidly, contraction arguments can replace the linear-in-$T$ envelope with tighter bounds.
Those refinements are treated in the spectral anonymity note (wiki) and can be imported here as needed.
\end{remark}

\begin{remark}[KL via Pinsker-type bounds]
If desired, TV bounds can be converted to KL using Pinsker-style inequalities (see, e.g., \cite{canonne_kl_tv}).
We avoid committing to a single conversion constant here; the exported interface is TV plus an explicit floor when ratio-style reasoning is required.
\end{remark}

\section{CoverSketch: a deployable pattern}
MC-EQ is a contract; CoverSketch is a pattern for realizing it.
The key insight is to choose a public cover support at each state and then equalize within that support.

\subsection{Public cover sets}
For each state $x\in S$, choose a public cover set $\mathcal{C}(x)\subseteq S$.
The reference kernel is typically
\[
K(y\mid x) = \begin{cases}1/|\mathcal{C}(x)| & y\in \mathcal{C}(x),\\ 0 & \text{otherwise.}\end{cases}
\]
This makes the floor parameter $\alpha = \min_x 1/|\mathcal{C}(x)|$ explicit.

\subsection{Equalization with slack}
A mechanism implements MC-EQ by ensuring its actual conditional next-state law stays within $\eps_t$ of this uniform reference.
Operationally, that slack is used to (i) respect feasibility constraints (committee availability), and (ii) adapt to transient performance goals.

\subsection{A minimal CoverSketch certificate}
For certification and release receipts \cite{series:release}, CoverSketch exports:
\begin{itemize}[leftmargin=1.2em]
\item the public state set $S$ and cover sets $\mathcal{C}(x)$;
\item the induced kernel $K$ and floor $\alpha$;
\item the per-step (or per-phase) slack schedule $(\eps_t)$;
\item any conditioning rules used to define $S$ (imported from \cite{series:state2}).
\end{itemize}
These are small, stable objects compared to raw traces.

\section{Design tradeoffs}
\paragraph{Sparsity vs. leakage.}
Smaller cover sets reduce cost but decrease $\alpha$ and therefore worsen $D_\infty$ bounds.
This makes the tradeoff explicit and auditable.

\paragraph{Mixing helps.}
If the reference chain $K$ mixes quickly, then even conservative TV union bounds can be improved using contraction arguments \cite{levin_peres_wilmer_mixing}.
In this series, those sharper bounds are developed in the spectral anonymity note (wiki).

\paragraph{Congestion coupling.}
Cover sets constrain routing choices and can interact with congestion.
We treat congestion observables separately via Congestion-EQ \cite{series:congeq}; the release workflow should export both certificates and treat them as linked budgets.

\section{Baselines and what they protect (discussion)}
\label{sec:baselines}
Table~\ref{tab:baseline-map} situates MC-EQ relative to common DHT privacy tactics.
The intent is not to survey the space (see the related-work routing map for that),\cite{series:relatedworkmap} but to make explicit \emph{which witness channel} each baseline addresses.

\begin{table}[t]
\centering
\small
\begin{tabular}{@{}p{0.28\linewidth}p{0.33\linewidth}p{0.31\linewidth}@{}}
\toprule
Approach & What it targets & Session risk / mismatch \\
\midrule
Anonymous-path DHT lookups (e.g., \textit{Octopus}) \cite{wang2012octopus} &
Initiator anonymity and robustness against active attacks &
Destination leakage can persist via history-dependent scheduling and retries unless equalization is explicit \\
OT/PIR-style query privacy \cite{backes2012queryprivacy,mazmudar2024peer2pir} &
Hides the queried key from intermediates/servers (crypto assumptions) &
Does not by itself equalize which committees/regions are contacted or control-plane patterns \\
Query obfuscation (double hashing / prefix fetching) \cite{daniel2025kademliaobf} &
Reduces interest leakage to individual responders &
Repetition may still amplify bias; conditioning on paths and retries can accumulate evidence \\
MUCC-style marginal equalization \cite{series:state2} &
Makes single-step contact distribution less distinguishing &
Not stable under history-conditioned schedulers without an explicit contract \\
\textbf{MC-EQ (this paper)} &
History-conditioned equalization to a public Markov reference &
Designed to remain meaningful under sessions; leakage scales with explicit $(T,k,\delta_{eq})$ (Table~\ref{tab:session-bound}) \\
\bottomrule
\end{tabular}
\caption{Baselines and how they relate to session adversaries.}
\label{tab:baseline-map}
\end{table}

\section{Deployment integration and receipt exports}
\label{sec:deploy}
MC-EQ/CoverSketch are designed to export verifiable line items for the release-receipt workflow.\cite{series:release}
A minimal deployment should export:
\begin{itemize}[leftmargin=1.3em]
\item equalization slack $\delta_{eq}$ (per step or per phase) and a reference-kernel identifier;
\item a support floor $q_{\min}$ (equivalently a cover-size bound $k_{\max}$ when using uniform covers);
\item the mixing/cover horizon $m$ and an explicit per-step contact budget $k$ (or committee fanout);
\item public lookup parameters that shape transcript length (e.g., concurrency/termination thresholds in iterative Kademlia-like deployments);
\item routing-table constraints that define the effective overlay graph (filters, admissible peers/committees); and
\item any retry/session-horizon assumptions used to translate per-step bounds into window budgets (Table~\ref{tab:session-bound}).
\end{itemize}

\paragraph{Exporting truncation vs.\ implementation deviation.}
CoverSketch typically targets a truncated row $\widetilde R_x$ (a top-$k$ or cover-restricted approximation of an ideal reference row $R_x$).
A conservative receipt should therefore separate
(i) a \emph{topology approximation} term $\eta_x:=\TV(\widetilde R_x,R_x)$ and
(ii) a \emph{mechanism deviation} term $\delta_{\mathrm{impl}}$ capturing how far the realized scheduler is from $\widetilde R_x$ under churn, packet loss, and implementation details.
By triangle inequality,
\[
\TV\big(P_d(\cdot\mid h_t),R_{x_t}\big)\le \delta_{\mathrm{impl}} + \eta_{x_t}.
\]
This keeps the bound honest without forcing the paper to model every deployment detail.

\paragraph{Churn and retries as explicit error terms.}
In practice, the realized transition kernel deviates from the intended reference due to churn (committee membership changes), packet loss, and retry logic.
A tight export is therefore an \emph{effective} per-step slack
$\delta_{eq}^{\mathrm{eff}}=\delta_{alg}+\delta_{churn}+\delta_{retry}$,
where $\delta_{churn}$ can be estimated by comparing observed transition frequencies to the advertised public topology over rolling windows.
Treating these as first-class receipt line items keeps the bound auditable.

\paragraph{Bundle record (optional).}
In the release workflow, MC-EQ/CoverSketch evidence can be summarized by an optional \texttt{mceq.json} record that commits to:
the public topology snapshot (e.g., a graph hash), the reference-kernel parameters, chosen $(m,k)$, per-state cover sets $\mathcal{C}(x)$, and tail masses $\eta_x$.
Downstream verifiers can replay the cover construction and check that the exported $(\delta_{eq},q_{\min})$ are realizable.

\paragraph{Engineering note.}
Even when destination equalization is correct, other control-plane surfaces may dominate leakage (indexers, gateways, provider advertisements, and monitoring beacons).
The release receipt should therefore treat MC-EQ as one line item among several and compose it with the companion tiers (retries/stop times, congestion witnesses, and control-plane releases).\cite{series:traceifaces,series:wcongeq,series:state3}

\section{Conclusion}
MC-EQ makes destination privacy modular: it is a local TV contract against a public reference chain.
CoverSketch is a simple, certificate-friendly pattern that instantiates that contract with explicit sparsity--leakage parameters.
Together, they let deployments publish destination privacy as a checkable release artifact, rather than an informal argument.

\section{Takeaways}
\begin{itemize}[leftmargin=*]
\item \textbf{Destination privacy needs its own witness:} even if queried keys are protected, \emph{where} traffic exits (destinations, providers, gateways) can be key-correlated without explicit accounting.
\item \textbf{MC-EQ/CoverSketch is a bridge:} treat destination leakage as an equalization problem over a destination sketch, so knobs (cover, mixing, committee choices) compile to a TV-style bound.
\item \textbf{Export line items for receipts:} the goal is a small manifest + accountant transcript that can be replayed and anchored in release receipts, not a narrative argument.
\end{itemize}

\begin{thebibliography}{99}\small
\bibitem{series:state2}
Anonymous.
\newblock \emph{AnonDHT State-Dependent Anonymity: Committees, Conditioning, and MUCC}.
\newblock AnonDHT state series Paper~2, The-One-True-Archive (2026).

\bibitem{series:release}
Anonymous.
\newblock \emph{Privacy as a Release Property: Proof-Carrying Anonymity Receipts for Anonymous DHTs}.
\newblock Release \& destination Paper~A, The-One-True-Archive (2026).


\bibitem{series:artifactpolicy}
Anonymous.
\newblock \emph{Artifact and Disclosure Policy for the One-True-Archive (Source-Only Public Release)}.
\newblock Series synthesis note (Synthesis~6), The-One-True-Archive (2026).

\bibitem{series:endpointbridge}
Anonymous.
\newblock \emph{Endpoint Metrics Bridge for Anonymous Overlays: Odds Inflation, Pointwise Maximal Leakage, and Maximal Leakage}.
\newblock Series synthesis note (Synthesis~5), The-One-True-Archive (2026).

\bibitem{series:congeq}
Anonymous.
\newblock \emph{Congestion-EQ: Total-Variation Budgets for Congestion Witnesses in Anonymous Overlays}.
\newblock Congestion series Paper~1, The-One-True-Archive (2026).

\bibitem{levin_peres_wilmer_mixing}
D.~A. Levin and Y.~Peres.
\newblock \emph{Markov Chains and Mixing Times (Second Edition)}.
\newblock AMS, 2017.

\bibitem{issa_maxleak}
I.~Issa, S.~Kamath, and A.~B. Wagner.
\newblock An operational approach to information leakage (maximal leakage).
\newblock arXiv:1807.07878, 2018.

\bibitem{canonne_kl_tv}
C.~L. Canonne.
\newblock A short note on an inequality between {KL} and {TV}.
\newblock arXiv:2202.07198, 2022.

\bibitem{series:recert}
Certified Series.
\emph{Disagreement-Gated Recertification}.
(Series paper in this archive: \texttt{certified\_series/paperC\_disagreement\_gated\_recertification/paper.tex}.)



\bibitem{series:schedcompiler}
Anonymous.
\newblock \emph{Scheduling and Compiler Techniques for Metadata-Hiding Overlay Lookups}.
\newblock Wiki: \texttt{[[2026.01.22 - Mathematics: Scheduling and Compiler Techniques for Metadata-Hiding Overlay Lookups]]}.

\bibitem{series:relatedworkmap}
Anonymous.
\newblock \emph{Related Work Map for Anonymity and Anonymous DHTs}.
\newblock Series synthesis note (Synthesis~7), The-One-True-Archive (2026).

\bibitem{series:traceifaces}
Anonymous.
\newblock \emph{Trace Interfaces for Anonymous DHTs: Tiers, Projections, and Witness Surfaces}.
\newblock Series synthesis note (Synthesis~3), The-One-True-Archive (2026).

\bibitem{series:wcongeq}
Anonymous.
\newblock \emph{W-Congestion-EQ: Proof-Carrying Congestion Budgets with Certificates and Transparency Logs}.
\newblock Congestion series Paper~3, The-One-True-Archive (2026).

\bibitem{series:state3}
Anonymous.
\newblock \emph{Closed-View Auditing and CPPC Manifests: Proof-Carrying Privacy Control Planes}.
\newblock AnonDHT state series Paper~3, The-One-True-Archive (2026).

\bibitem{wang2012octopus}
Wang et al.
\newblock \emph{Octopus: A secure and anonymous DHT lookup}.
\newblock In \emph{Proceedings of NDSS}, 2012.

\bibitem{backes2012queryprivacy}
Backes et al.
\newblock \emph{Adding query privacy to robust DHTs}.
\newblock In \emph{Proceedings of ESORICS}, 2012.

\bibitem{mazmudar2024peer2pir}
Mazmudar et al.
\newblock \emph{Peer2PIR: Practical private information retrieval for peer-to-peer content routing}.
\newblock Preprint / technical report, 2024.

\bibitem{daniel2025kademliaobf}
Daniel et al.
\newblock \emph{Query obfuscation methods for Kademlia-style lookups}.
\newblock Preprint / draft, 2025.



\bibitem{series:notation}
Anonymous.
\newblock Notation and Minimal Model for the One-True-Archive.
\newblock Series synthesis note (Synthesis~8), 2026. (This archive.)

\end{thebibliography}

\end{document}